%% ============================================================
%% ebook.latex — Pandoc LaTeX template for Code4Projects ebooks
%% ============================================================
\documentclass[11pt,american,a4paper,openany]{book}

%% ── Encoding & language ─────────────────────────────────────
\usepackage[utf8]{inputenc}
\usepackage[T1]{fontenc}
\usepackage[shorthands=off,main=american]{babel}

%% ── Fonts (xelatex) ─────────────────────────────────────────
\usepackage{fontspec}
\setmainfont{Georgia}
\setsansfont{Helvetica Neue}
\setmonofont[Scale=0.88]{Menlo}

%% ── Page geometry ───────────────────────────────────────────
\usepackage[
  left=2.5cm,
  right=2.5cm,
  top=3cm,
  bottom=3cm,
]{geometry}

%% ── Line spacing ────────────────────────────────────────────
\usepackage{setspace}
\setstretch{1.4}

%% ── Colors & links ──────────────────────────────────────────
\usepackage[dvipsnames,svgnames,x11names]{xcolor}
\usepackage[
  colorlinks=true,
  linkcolor=NavyBlue,
  urlcolor=NavyBlue,
  citecolor=NavyBlue,
  filecolor=NavyBlue,
  pdfusetitle
]{hyperref}

%% ── Code listings ───────────────────────────────────────────
\usepackage{fancyvrb}
\usepackage{fvextra}

\definecolor{codebg}{HTML}{1e1e2e}
\definecolor{codefg}{HTML}{cdd6f4}
\definecolor{codeframe}{HTML}{313244}

\RecustomVerbatimEnvironment{verbatim}{Verbatim}{
  fontsize=\small,
  frame=single,
  framesep=6pt,
  rulecolor=\color{codeframe},
  fillcolor=\color{codebg},
  formatcom=\color{codefg},
  breaklines=true,
  breakanywhere=true,
  xleftmargin=10pt,
  xrightmargin=10pt,
}

%% ── Tables ──────────────────────────────────────────────────
\usepackage{longtable}
\usepackage{booktabs}
\usepackage{array}

%% ── Images ──────────────────────────────────────────────────
\usepackage{graphicx}
\makeatletter
\def\maxwidth{\ifdim\Gin@nat@width>\linewidth\linewidth\else\Gin@nat@width\fi}
\def\maxheight{\ifdim\Gin@nat@height>\textheight\textheight\else\Gin@nat@height\fi}
\makeatother
\setkeys{Gin}{width=\maxwidth,height=\maxheight,keepaspectratio}

%% ── Misc ────────────────────────────────────────────────────
\usepackage{parskip}
\setlength{\parindent}{0pt}
\usepackage{microtype}
\usepackage{enumitem}
\usepackage{titlesec}
\usepackage{mdframed}

%% Chapter style
\titleformat{\chapter}[display]
  {\normalfont\huge\bfseries\sffamily}
  {\chaptertitlename\ \thechapter}{16pt}{\Huge}
\titlespacing*{\chapter}{0pt}{30pt}{20pt}

%% Section style
\titleformat{\section}
  {\normalfont\Large\bfseries\sffamily}
  {\thesection}{1em}{}
\titleformat{\subsection}
  {\normalfont\large\bfseries\sffamily}
  {\thesubsection}{1em}{}

%% ── Headers / footers ───────────────────────────────────────
\usepackage{fancyhdr}
\pagestyle{fancy}
\fancyhf{}
\fancyhead[LE,RO]{\small\thepage}
\fancyhead[RE]{\small\itshape Getting Started with Docker}
\fancyhead[LO]{\small\itshape\leftmark}
\renewcommand{\headrulewidth}{0.4pt}

%% ── Blockquotes ─────────────────────────────────────────────
\newmdenv[
  leftline=true, rightline=false, topline=false, bottomline=false,
  linewidth=3pt, linecolor=NavyBlue!60,
  backgroundcolor=NavyBlue!5,
  innerleftmargin=12pt, innerrightmargin=12pt,
  innertopmargin=8pt, innerbottommargin=8pt,
  skipabove=10pt, skipbelow=10pt,
]{quoteblock}
\renewenvironment{quote}{\begin{quoteblock}}{\end{quoteblock}}

%% ── pandoc required macros ──────────────────────────────────
\providecommand{\tightlist}{%
  \setlength{\itemsep}{2pt}\setlength{\parskip}{0pt}}

\usepackage{color}
\usepackage{fancyvrb}
\newcommand{\VerbBar}{|}
\newcommand{\VERB}{\Verb[commandchars=\\\{\}]}
\DefineVerbatimEnvironment{Highlighting}{Verbatim}{commandchars=\\\{\}}
% Add ',fontsize=\small' for more characters per line
\newenvironment{Shaded}{}{}
\newcommand{\AlertTok}[1]{\textcolor[rgb]{1.00,0.00,0.00}{\textbf{#1}}}
\newcommand{\AnnotationTok}[1]{\textcolor[rgb]{0.38,0.63,0.69}{\textbf{\textit{#1}}}}
\newcommand{\AttributeTok}[1]{\textcolor[rgb]{0.49,0.56,0.16}{#1}}
\newcommand{\BaseNTok}[1]{\textcolor[rgb]{0.25,0.63,0.44}{#1}}
\newcommand{\BuiltInTok}[1]{\textcolor[rgb]{0.00,0.50,0.00}{#1}}
\newcommand{\CharTok}[1]{\textcolor[rgb]{0.25,0.44,0.63}{#1}}
\newcommand{\CommentTok}[1]{\textcolor[rgb]{0.38,0.63,0.69}{\textit{#1}}}
\newcommand{\CommentVarTok}[1]{\textcolor[rgb]{0.38,0.63,0.69}{\textbf{\textit{#1}}}}
\newcommand{\ConstantTok}[1]{\textcolor[rgb]{0.53,0.00,0.00}{#1}}
\newcommand{\ControlFlowTok}[1]{\textcolor[rgb]{0.00,0.44,0.13}{\textbf{#1}}}
\newcommand{\DataTypeTok}[1]{\textcolor[rgb]{0.56,0.13,0.00}{#1}}
\newcommand{\DecValTok}[1]{\textcolor[rgb]{0.25,0.63,0.44}{#1}}
\newcommand{\DocumentationTok}[1]{\textcolor[rgb]{0.73,0.13,0.13}{\textit{#1}}}
\newcommand{\ErrorTok}[1]{\textcolor[rgb]{1.00,0.00,0.00}{\textbf{#1}}}
\newcommand{\ExtensionTok}[1]{#1}
\newcommand{\FloatTok}[1]{\textcolor[rgb]{0.25,0.63,0.44}{#1}}
\newcommand{\FunctionTok}[1]{\textcolor[rgb]{0.02,0.16,0.49}{#1}}
\newcommand{\ImportTok}[1]{\textcolor[rgb]{0.00,0.50,0.00}{\textbf{#1}}}
\newcommand{\InformationTok}[1]{\textcolor[rgb]{0.38,0.63,0.69}{\textbf{\textit{#1}}}}
\newcommand{\KeywordTok}[1]{\textcolor[rgb]{0.00,0.44,0.13}{\textbf{#1}}}
\newcommand{\NormalTok}[1]{#1}
\newcommand{\OperatorTok}[1]{\textcolor[rgb]{0.40,0.40,0.40}{#1}}
\newcommand{\OtherTok}[1]{\textcolor[rgb]{0.00,0.44,0.13}{#1}}
\newcommand{\PreprocessorTok}[1]{\textcolor[rgb]{0.74,0.48,0.00}{#1}}
\newcommand{\RegionMarkerTok}[1]{#1}
\newcommand{\SpecialCharTok}[1]{\textcolor[rgb]{0.25,0.44,0.63}{#1}}
\newcommand{\SpecialStringTok}[1]{\textcolor[rgb]{0.73,0.40,0.53}{#1}}
\newcommand{\StringTok}[1]{\textcolor[rgb]{0.25,0.44,0.63}{#1}}
\newcommand{\VariableTok}[1]{\textcolor[rgb]{0.10,0.09,0.49}{#1}}
\newcommand{\VerbatimStringTok}[1]{\textcolor[rgb]{0.25,0.44,0.63}{#1}}
\newcommand{\WarningTok}[1]{\textcolor[rgb]{0.38,0.63,0.69}{\textbf{\textit{#1}}}}

%% ============================================================
%% DOCUMENT
%% ============================================================
\begin{document}

%% ── Title page ──────────────────────────────────────────────
\pagestyle{empty}
\vspace*{2cm}
\begin{center}
  {\fontsize{36}{44}\selectfont\bfseries\sffamily Getting Started with
Docker}\\[0.8em]
  {\Large\sffamily\color{gray} Containers, Images, Networking, Volumes,
Compose and Security}\\[3em]
  {\large Salvatore D'Angelo}\\[0.5em]
  {\normalsize\color{gray} 2025}\\[4em]
  {\small\color{gray} © 2025 Salvatore D'Angelo --- code4projects.com.
All rights reserved.}
\end{center}
\newpage

%% ── Table of contents ───────────────────────────────────────
\pagestyle{plain}
{
\hypersetup{linkcolor=NavyBlue}
\setcounter{tocdepth}{2}
\tableofcontents
}
\newpage

%% ── Body ────────────────────────────────────────────────────
\pagestyle{fancy}
\section{Getting Started with Docker and
Podman}\label{getting-started-with-docker-and-podman}

\emph{Posted on ****}

\subsection{Introduction}\label{introduction}

This is the first article of the \textbf{Getting Started with Docker}
series. By the end of this article you will understand what Docker is,
what problems it solves, and you will have your first containerized web
server running on your machine in minutes --- without writing a single
line of Dockerfile.

You should read this article if:

\begin{itemize}
\tightlist
\item
  You have heard of Docker but never used it
\item
  You want a practical, no-nonsense introduction that gets you running
  immediately
\item
  You are on a machine where Docker Desktop is not available and you
  need an alternative (Podman)
\end{itemize}

All the commands shown here work identically with both \textbf{Docker}
and \textbf{Podman}. Where there are differences, they are explicitly
noted.

\subsection{Why Docker?}\label{why-docker}

Before diving into commands, it is worth spending a minute on the
problem Docker actually solves.

Imagine you are a developer. You write an application on your laptop, it
works perfectly. Then you hand it off to a colleague or deploy it to a
server --- and it breaks. Different OS version, different library
installed, different path. This is the classic \textbf{``it works on my
machine''} problem, and it has tormented developers for decades.

Docker solves this by packaging your application together with
everything it needs to run --- runtime, libraries, configuration ---
into a single portable unit called a \textbf{container}. That container
runs identically on your laptop, your colleague's machine, and the
production server.

This matters not just for individual developers. Modern applications are
built as dozens (sometimes hundreds) of independent services ---
\textbf{microservices} --- each with its own dependencies. Docker makes
it possible to deploy, update, and scale each service independently,
without interference.

\subsection{Docker vs Podman: Which One Should You
Use?}\label{docker-vs-podman-which-one-should-you-use}

Both Docker and Podman are tools for running containers. For everything
covered in this series, they are interchangeable --- the CLI commands
are identical.

The key differences:

\begin{longtable}[]{@{}
  >{\raggedright\arraybackslash}p{(\linewidth - 4\tabcolsep) * \real{0.3333}}
  >{\raggedright\arraybackslash}p{(\linewidth - 4\tabcolsep) * \real{0.3333}}
  >{\raggedright\arraybackslash}p{(\linewidth - 4\tabcolsep) * \real{0.3333}}@{}}
\toprule\noalign{}
\begin{minipage}[b]{\linewidth}\raggedright
\end{minipage} & \begin{minipage}[b]{\linewidth}\raggedright
Docker
\end{minipage} & \begin{minipage}[b]{\linewidth}\raggedright
Podman
\end{minipage} \\
\midrule\noalign{}
\endhead
\bottomrule\noalign{}
\endlastfoot
\textbf{Architecture} & Requires a background daemon (\texttt{dockerd})
& Daemonless --- runs containers directly \\
\textbf{Privileges} & Daemon runs as root & Rootless by default --- more
secure \\
\textbf{License} & Docker Desktop requires a paid license for enterprise
use & Fully open source (Apache 2.0) \\
\textbf{Installation on Mac} & Docker Desktop &
\texttt{brew\ install\ podman} + \texttt{podman\ machine\ init} \\
\textbf{Compose support} & \texttt{docker\ compose} (v2 plugin) &
\texttt{podman\ compose} (from Podman 4.7+) \\
\textbf{OCI compatibility} & Yes & Yes --- same image format \\
\end{longtable}

\textbf{In short}: if Docker Desktop is available and licensed for your
use, use Docker. If you are on a corporate machine where Docker Desktop
requires a commercial license, or if you prefer an open source rootless
alternative, use Podman. The rest of this series works exactly the same
either way.

\begin{quote}
Throughout this series, all commands are shown with \texttt{docker}. If
you are using Podman, simply replace \texttt{docker} with
\texttt{podman} --- or set the alias \texttt{alias\ docker=podman} once
and forget about it.
\end{quote}

\subsection{Installing Docker or
Podman}\label{installing-docker-or-podman}

\subsubsection{Docker}\label{docker}

On Linux, install Docker Engine following the
\href{https://docs.docker.com/engine/install/}{official documentation}.
On Mac and Windows, Docker Desktop is the easiest option if your license
allows it: download it from
\href{https://docs.docker.com/desktop/}{docs.docker.com/desktop}.

\subsubsection{Podman (Mac)}\label{podman-mac}

On macOS, Podman uses a lightweight Linux virtual machine behind the
scenes --- the same approach Docker Desktop uses. Install and initialise
it with three commands:

\begin{Shaded}
\begin{Highlighting}[]
\NormalTok{brew install podman}
\NormalTok{podman machine init}
\NormalTok{podman machine start}
\end{Highlighting}
\end{Shaded}

After this, \texttt{podman} is ready and behaves exactly like
\texttt{docker}.

\subsection{Core Concepts: Images and
Containers}\label{core-concepts-images-and-containers}

Before running anything, two concepts are essential.

\subsubsection{Docker Image}\label{docker-image}

A \textbf{Docker image} is a read-only, layered package that contains
everything needed to run a specific application: the OS filesystem,
runtime, libraries, configuration files, and the application itself.

Images are stored in registries. The most popular is
\href{https://hub.docker.com}{Docker Hub}, which hosts thousands of
official and community images. When you run a container, Docker (or
Podman) downloads the image from the registry if it is not already
present on your machine.

Images are built in \textbf{layers}. Each instruction that modifies the
filesystem adds a layer on top of the previous ones. Layers are cached
and shared between images --- this is what keeps images small and builds
fast.

\subsubsection{Docker Container}\label{docker-container}

A \textbf{container} is a running instance of an image. You can think of
the image as a class and the container as an object instantiated from
it. Multiple containers can run from the same image simultaneously, each
fully isolated from the others.

Containers share the host OS kernel but are isolated from each other
using Linux kernel features: \textbf{namespaces} (for process, network,
filesystem isolation) and \textbf{cgroups} (for resource limits like CPU
and memory).

This is the fundamental difference from a Virtual Machine: a VM emulates
an entire hardware stack and runs a separate OS kernel. A container
shares the host kernel and is much lighter as a result.

We will explore this topic in depth in the
\href{/Users/sasadangelo/github.com/sasadangelo/code4projects/containers-vs-virtual-machines-2025/}{next
article}.

\subsection{Essential Commands}\label{essential-commands}

Here is the minimal set of commands you need to start working with
Docker. These are the ones you will use every day.

\subsubsection{List local images}\label{list-local-images}

\begin{Shaded}
\begin{Highlighting}[]
\NormalTok{docker image ls}
\end{Highlighting}
\end{Shaded}

At the beginning, the list is empty.

\subsubsection{Pull an image from Docker
Hub}\label{pull-an-image-from-docker-hub}

\begin{Shaded}
\begin{Highlighting}[]
\NormalTok{docker pull nginx:alpine}
\end{Highlighting}
\end{Shaded}

This downloads the official Nginx image based on Alpine Linux --- a very
small (\textasciitilde40 MB) but fully functional web server image.

\subsubsection{Run a container}\label{run-a-container}

\begin{Shaded}
\begin{Highlighting}[]
\NormalTok{docker run {-}d {-}p 8080:80 {-}{-}name my{-}nginx nginx:alpine}
\end{Highlighting}
\end{Shaded}

Breaking down the options:

\begin{itemize}
\tightlist
\item
  \texttt{-d} --- runs the container in detached mode (in the
  background, shell does not hang)
\item
  \texttt{-p\ 8080:80} --- maps port 8080 on your host to port 80 inside
  the container
\item
  \texttt{-\/-name\ my-nginx} --- gives the container a friendly name
  instead of a random one
\item
  \texttt{nginx:alpine} --- the image to use
\end{itemize}

Open your browser at \texttt{http://localhost:8080} and you will see the
Nginx welcome page. You just ran a web server without installing
anything on your machine.

\subsubsection{List running containers}\label{list-running-containers}

\begin{Shaded}
\begin{Highlighting}[]
\NormalTok{docker container ls}
\end{Highlighting}
\end{Shaded}

To see all containers including stopped ones:

\begin{Shaded}
\begin{Highlighting}[]
\NormalTok{docker container ls {-}a}
\end{Highlighting}
\end{Shaded}

\subsubsection{View container logs}\label{view-container-logs}

\begin{Shaded}
\begin{Highlighting}[]
\NormalTok{docker logs my{-}nginx}
\end{Highlighting}
\end{Shaded}

\subsubsection{Execute a command inside a running
container}\label{execute-a-command-inside-a-running-container}

\begin{Shaded}
\begin{Highlighting}[]
\NormalTok{docker exec {-}it my{-}nginx /bin/sh}
\end{Highlighting}
\end{Shaded}

The \texttt{-it} flags attach an interactive terminal. Note:
Alpine-based images use \texttt{sh} instead of \texttt{bash}.

\subsubsection{Stop and remove a
container}\label{stop-and-remove-a-container}

\begin{Shaded}
\begin{Highlighting}[]
\NormalTok{docker stop my{-}nginx}
\NormalTok{docker rm my{-}nginx}
\end{Highlighting}
\end{Shaded}

Or in one command:

\begin{Shaded}
\begin{Highlighting}[]
\NormalTok{docker rm {-}f my{-}nginx}
\end{Highlighting}
\end{Shaded}

\subsubsection{Remove an image}\label{remove-an-image}

\begin{Shaded}
\begin{Highlighting}[]
\NormalTok{docker image rm nginx:alpine}
\end{Highlighting}
\end{Shaded}

\subsubsection{Clean up everything
unused}\label{clean-up-everything-unused}

\begin{Shaded}
\begin{Highlighting}[]
\NormalTok{docker system prune}
\end{Highlighting}
\end{Shaded}

This removes all stopped containers, dangling images, and unused
networks. Add \texttt{-a} to also remove unused images that are not
referenced by any container.

\subsection{Your First Containerized
Application}\label{your-first-containerized-application}

Let us put it all together with a practical example. You will serve a
custom HTML page through Nginx --- without installing Nginx, without a
Dockerfile, using only the official image and a single
\texttt{docker\ run} command.

Create a file called \texttt{index.html} anywhere on your machine:

\begin{Shaded}
\begin{Highlighting}[]

\NormalTok{  Hello from Docker}
  
\NormalTok{    Hello, World!}
\NormalTok{    This page is served by Nginx running inside a container.}
  
\end{Highlighting}
\end{Shaded}

Now run Nginx and mount that file into the container:

\begin{Shaded}
\begin{Highlighting}[]
\NormalTok{docker run {-}d \textbackslash{}}
\NormalTok{  {-}p 8080:80 \textbackslash{}}
\NormalTok{  {-}{-}name hello{-}nginx \textbackslash{}}
\NormalTok{  {-}v $(pwd)/index.html:/usr/share/nginx/html/index.html:ro \textbackslash{}}
\NormalTok{  nginx:alpine}
\end{Highlighting}
\end{Shaded}

The \texttt{-v} option \textbf{bind mounts} your local file into the
container's web root. The \texttt{:ro} suffix makes the mount read-only
inside the container (a good practice).

Open \texttt{http://localhost:8080} in your browser:

\begin{quote}
\textbf{Hello, World!} This page is served by Nginx running inside a
container.
\end{quote}

You served a custom page without installing a web server, without
writing a Dockerfile, and without touching your system configuration.
When you are done:

\begin{Shaded}
\begin{Highlighting}[]
\NormalTok{docker rm {-}f hello{-}nginx}
\end{Highlighting}
\end{Shaded}

Your machine is exactly as it was before. Nothing installed, nothing
left behind.

\subsection{How This Series Is
Structured}\label{how-this-series-is-structured}

This first article introduced the core concepts and got you running
immediately. In the upcoming articles we will progressively build on
this foundation:

\begin{itemize}
\tightlist
\item
  \textbf{Containers vs Virtual Machines} --- deep dive into isolation
  mechanisms, namespaces, cgroups, and the difference from
  hypervisor-based virtualization
\item
  \textbf{Dockerfile \& Building Custom Images} --- write your own
  \texttt{Dockerfile}, build and tag a custom image, understand the
  layer cache
\item
  \textbf{How Docker Networking Works} --- user-defined bridge networks,
  automatic DNS resolution between containers, network isolation
\item
  \textbf{How Docker Volumes Work} --- named volumes vs bind mounts,
  persisting data across container restarts
\item
  \textbf{How Docker Compose Works} --- manage multi-container
  applications with a single \texttt{docker\ compose\ up}
\item
  \textbf{Docker Security Best Practices} --- run containers as
  non-root, drop capabilities, scan images for vulnerabilities
\end{itemize}

Each article builds on the previous one around the same practical Nginx
application, adding one concept at a time.

\subsection{Conclusion}\label{conclusion}

In this article we covered:

\begin{itemize}
\tightlist
\item
  The problem Docker solves and why containers matter
\item
  The difference between Docker and Podman, and when to use each
\item
  The two fundamental concepts: \textbf{images} and \textbf{containers}
\item
  The essential CLI commands for daily use
\item
  Running a custom web page with Nginx using only the official image
\end{itemize}

The
\href{/Users/sasadangelo/github.com/sasadangelo/code4projects/containers-vs-virtual-machines-2025/}{next
article} goes deeper into the conceptual foundations: what actually
makes a container different from a virtual machine, and what Linux
kernel mechanisms make it all work.

\begin{center}\rule{0.5\linewidth}{0.5pt}\end{center}

If you enjoyed this article, don't forget to \textbf{give it a clap 👏},
\textbf{share it with your friends 🔗}, and \textbf{follow me for more
tips and tutorials on software development 📘}. Your support helps me
create more content like this --- thank you! 🙌

\section{Containers vs Virtual
Machines}\label{containers-vs-virtual-machines}

\emph{Posted on ****}

\subsection{Introduction}\label{introduction-1}

This is the second article of the \textbf{Getting Started with Docker}
series. In the
\href{/Users/sasadangelo/github.com/sasadangelo/code4projects/getting-started-with-docker-2025/}{previous
article} we introduced Docker and ran our first container. Before going
further, it is worth spending time on a fundamental question:
\textbf{what is the real difference between a container and a virtual
machine?}

When I started working with Docker this was not easy to grasp. The two
technologies look similar from the outside --- both give you an isolated
environment to run software --- but they work in completely different
ways under the hood.

Virtualization can be achieved in two ways:

\begin{itemize}
\tightlist
\item
  \textbf{Full virtualization} --- a hypervisor simulates enough
  hardware to run a complete guest operating system. This is the Virtual
  Machine approach.
\item
  \textbf{Operating system-level virtualization} --- the OS kernel
  itself is used to create isolated environments that share the same
  kernel. This is the Container approach.
\end{itemize}

\subsection{Virtual Machines}\label{virtual-machines}

Historically, as server processing power grew, applications running on
bare metal couldn't exploit the abundance of available resources.
\textbf{Virtual Machines} were born to solve this: by running a
hypervisor on top of physical hardware, you could carve up one server
into many independent guest systems.

A \textbf{hypervisor} (also called Virtual Machine Monitor) is the
software layer that sits between the hardware and the virtual machines.
It can be type 1 (bare metal, like VMware ESXi or KVM) or type 2
(hosted, like VirtualBox or Parallels).

Each VM includes:

\begin{itemize}
\tightlist
\item
  A full copy of a guest operating system
\item
  Virtual hardware (CPU, RAM, disk, network)
\item
  All the application dependencies
\end{itemize}

The system running the VMs is called the \textbf{host}, the systems
running inside are called \textbf{guests}.

\subsubsection{Pros and cons of VMs}\label{pros-and-cons-of-vms}

\begin{verbatim}
    Pros
  
  
    Strong isolation — each VM has its own kernel.
    Run different operating systems on the same host (Windows on Linux, etc.).
    Create and destroy on demand in minutes.
    Snapshot a VM at any point in time.
    Foundation of the IaaS cloud model (AWS EC2, Azure VMs, etc.).
  


    Cons
  
  
    Heavy — each VM carries a full OS (GBs of disk, hundreds of MB of RAM).
    Slow to boot — minutes, not seconds.
    Hypervisor adds overhead to application execution.
    Managing a large fleet of VMs is operationally expensive.
  
\end{verbatim}

The rise of VMs gave birth to the \textbf{Infrastructure as a Service
(IaaS)} market. Companies like Amazon, Microsoft, and IBM allow you to
rent virtual machines on demand. A startup today does not need to buy
physical hardware --- it registers an account with a cloud provider,
orders a VM with the required specs, and is ready in minutes.

\subsection{Containers}\label{containers}

A \textbf{container} is an isolated environment created using features
of the Linux kernel itself --- not a separate hypervisor. Multiple
containers run on the same host, all sharing the same OS kernel, but
each seeing its own isolated filesystem, processes, network, and users.

The two kernel mechanisms that make this possible are:

\begin{itemize}
\tightlist
\item
  \textbf{Namespaces} --- give each container an isolated view of the
  system: its own process tree (PID namespace), network interfaces (net
  namespace), filesystem (mount namespace), users (user namespace), and
  more.
\item
  \textbf{cgroups} (control groups) --- limit and account for resource
  usage (CPU, memory, disk I/O) per container.
\end{itemize}

Together, namespaces and cgroups create the illusion of a separate
machine without the overhead of emulating hardware or running a second
kernel.

\begin{quote}
\textbf{Important note for Mac and Windows users}: because containers
rely on Linux kernel features, Docker and Podman on macOS and Windows
run a lightweight Linux VM behind the scenes (using Apple's
Virtualization Framework on Mac, and WSL2 on Windows). Your containers
still run in that Linux environment --- you just don't see it.
\end{quote}

\subsubsection{Pros and cons of
containers}\label{pros-and-cons-of-containers}

\begin{verbatim}
    Pros
  
  
    Lightweight — a container image is tens or hundreds of MB, not GBs.
    Fast — containers start in milliseconds to seconds.
    No overhead from emulating hardware.
    Portable — same image runs on your laptop, CI server, and production.
    Solves the "it works on my machine" problem.
    Natural fit for microservices architecture.
  


    Cons
  
  
    Weaker isolation than VMs — all containers share the host kernel.
    Cannot run a different OS kernel (e.g., no Windows containers on Linux).
    Managing many containers requires an orchestrator (Kubernetes, Swarm).
    Security requires care — a misconfigured container can affect the host.
  
\end{verbatim}

\subsection{The Analogy}\label{the-analogy}

Mike Coleman on the
\href{https://blog.docker.com/2016/03/containers-are-not-vms/}{Docker
official blog} provides a memorable analogy:

\begin{quote}
Houses (VMs) are fully self-contained and offer protection from unwanted
guests. They each possess their own infrastructure --- plumbing,
heating, electrical. Apartments (containers) also offer protection from
unwanted guests, but they are built around shared infrastructure. The
apartment building (Docker Host) shares plumbing, heating, electrical.
Apartments are offered in all kinds of sizes --- from studio to
penthouse. You only rent exactly what you need.
\end{quote}

\subsection{Containers and VMs: not enemies,
complementaries}\label{containers-and-vms-not-enemies-complementaries}

In practice, containers and VMs are often used \textbf{together}, not in
opposition. The most common production setup is:

\begin{enumerate}
\def\labelenumi{\arabic{enumi}.}
\tightlist
\item
  A cloud provider gives you a VM (an EC2 instance, an Azure VM, etc.)
\item
  Docker or a container runtime runs inside that VM
\item
  Your application runs as containers inside Docker
\end{enumerate}

VMs provide the strong isolation boundary at the infrastructure level.
Containers provide the lightweight, fast, portable packaging for
applications within that boundary.

\subsection{Conclusion}\label{conclusion-1}

In this article we covered:

\begin{itemize}
\tightlist
\item
  The difference between full virtualization (VMs) and OS-level
  virtualization (containers)
\item
  How \textbf{namespaces} and \textbf{cgroups} implement container
  isolation
\item
  Pros and cons of each technology
\item
  Why containers and VMs are complementary, not competing
\end{itemize}

The
\href{/Users/sasadangelo/github.com/sasadangelo/code4projects/dockerfile-and-building-custom-images/}{next
article} focuses on \textbf{Dockerfiles}: how to write your own image
definition and build a custom container image from scratch.

\begin{center}\rule{0.5\linewidth}{0.5pt}\end{center}

If you enjoyed this article, don't forget to \textbf{give it a clap 👏},
\textbf{share it with your friends 🔗}, and \textbf{follow me for more
tips and tutorials on software development 📘}. Your support helps me
create more content like this --- thank you! 🙌

\section{Dockerfile and Building Custom
Images}\label{dockerfile-and-building-custom-images}

\emph{Posted on ****}

\subsection{Introduction}\label{introduction-2}

This is the third article of the \textbf{Getting Started with Docker}
series. In the
\href{/Users/sasadangelo/github.com/sasadangelo/code4projects/containers-vs-virtual-machines-2025/}{previous
article} we explored the differences between containers and virtual
machines. In the
\href{/Users/sasadangelo/github.com/sasadangelo/code4projects/getting-started-with-docker-2025/}{first
article} we ran Nginx directly from the official image using a bind
mount to serve a custom page.

That approach works well for quick experiments, but in a real project
you want a \textbf{reproducible, self-contained image} that already
includes your content and configuration --- one that anyone on your team
(or any CI/CD pipeline) can build and run without manual steps.

This is where \textbf{Dockerfiles} come in.

\subsection{What is a Dockerfile?}\label{what-is-a-dockerfile}

A \textbf{Dockerfile} is a plain-text file containing a sequence of
instructions that tell Docker how to build an image. Each instruction
creates a new \textbf{layer} on top of the previous one. When you
rebuild an image, Docker only rebuilds the layers that changed ---
everything before the first change is served from the \textbf{layer
cache}, making subsequent builds very fast.

\subsection{Dockerfile Instructions}\label{dockerfile-instructions}

Here are the most important instructions you will use in every
Dockerfile.

\subsubsection{FROM}\label{from}

\begin{Shaded}
\begin{Highlighting}[]
\KeywordTok{FROM}\NormalTok{ nginx:alpine}
\end{Highlighting}
\end{Shaded}

Every Dockerfile must start with \texttt{FROM}. It specifies the
\textbf{base image} --- the starting point. In this series we always
start from \texttt{nginx:alpine}, a minimal image that already contains
a working Nginx web server.

\subsubsection{COPY}\label{copy}

\begin{Shaded}
\begin{Highlighting}[]
\KeywordTok{COPY}\NormalTok{ ./html /usr/share/nginx/html}
\end{Highlighting}
\end{Shaded}

Copies files or directories from your build context (the folder where
you run \texttt{docker\ build}) into the image. This is the main
instruction for adding your application code or configuration.

\subsubsection{ADD}\label{add}

\texttt{ADD} is similar to \texttt{COPY} but with two extra powers: it
can automatically unpack compressed archives (\texttt{.tar.gz}, etc.)
and it can fetch files from remote URLs. For simply copying local files,
\textbf{prefer \texttt{COPY}} --- it is explicit and predictable.

\subsubsection{RUN}\label{run}

\begin{Shaded}
\begin{Highlighting}[]
\KeywordTok{RUN} \ExtensionTok{apk}\NormalTok{ add }\AttributeTok{{-}{-}no{-}cache}\NormalTok{ curl}
\end{Highlighting}
\end{Shaded}

Executes a command \textbf{during the build} and commits the result as a
new layer. Use it to install packages, compile code, or run any setup
step. Each \texttt{RUN} instruction is a layer, so chain related
commands with \texttt{\&\&} to minimize layer count.

\subsubsection{ENV}\label{env}

\begin{Shaded}
\begin{Highlighting}[]
\KeywordTok{ENV}\NormalTok{ APP\_ENV=production}
\end{Highlighting}
\end{Shaded}

Sets an environment variable that is available both during the build and
at container runtime.

\subsubsection{EXPOSE}\label{expose}

\begin{Shaded}
\begin{Highlighting}[]
\KeywordTok{EXPOSE}\NormalTok{ 80}
\end{Highlighting}
\end{Shaded}

Documents which port the container listens on. It does \textbf{not}
actually publish the port --- that happens at runtime with \texttt{-p}.
Think of it as metadata for the image user.

\subsubsection{CMD}\label{cmd}

\begin{Shaded}
\begin{Highlighting}[]
\KeywordTok{CMD}\NormalTok{ [}\StringTok{"nginx"}\NormalTok{, }\StringTok{"{-}g"}\NormalTok{, }\StringTok{"daemon off;"}\NormalTok{]}
\end{Highlighting}
\end{Shaded}

Specifies the default command to run when a container starts. It can be
overridden at runtime. Use the \textbf{exec form} (JSON array) rather
than the shell form to avoid wrapping the process in a shell --- this
ensures signals like \texttt{SIGTERM} are delivered directly to your
process.

\subsubsection{ENTRYPOINT}\label{entrypoint}

Similar to \texttt{CMD} but not easily overridden. Use
\texttt{ENTRYPOINT} when the container has a fixed main executable, and
\texttt{CMD} to provide default arguments to it.

\subsubsection{USER}\label{user}

\begin{Shaded}
\begin{Highlighting}[]
\KeywordTok{USER}\NormalTok{ nginx}
\end{Highlighting}
\end{Shaded}

Switches to a non-root user for all subsequent instructions and at
runtime. This is a \textbf{security best practice} --- we will cover it
in detail in the security article.

\subsection{Building a Custom Nginx
Image}\label{building-a-custom-nginx-image}

Let us build a custom image that already includes our HTML page, so
anyone can run it with a single \texttt{docker\ run} --- no bind mounts,
no extra files needed.

Create a project directory:

\begin{Shaded}
\begin{Highlighting}[]
\NormalTok{mkdir my{-}nginx \&\& cd my{-}nginx}
\end{Highlighting}
\end{Shaded}

Create the HTML content:

\begin{Shaded}
\begin{Highlighting}[]
\NormalTok{mkdir html}
\NormalTok{cat \textgreater{} html/index.html }

\NormalTok{  My Custom Nginx}
  
\NormalTok{    Hello from my custom image!}
\NormalTok{    This page is baked into the Docker image.}
  

\NormalTok{EOF}
\end{Highlighting}
\end{Shaded}

Create the Dockerfile:

\begin{Shaded}
\begin{Highlighting}[]
\KeywordTok{FROM}\NormalTok{ nginx:alpine}

\CommentTok{\# Copy our HTML content into the image}
\KeywordTok{COPY}\NormalTok{ html /usr/share/nginx/html}

\CommentTok{\# Document which port Nginx listens on}
\KeywordTok{EXPOSE}\NormalTok{ 80}
\end{Highlighting}
\end{Shaded}

Build the image:

\begin{Shaded}
\begin{Highlighting}[]
\NormalTok{docker build {-}t my{-}nginx:1.0 .}
\end{Highlighting}
\end{Shaded}

Breaking down the command:

\begin{itemize}
\tightlist
\item
  \texttt{-t\ my-nginx:1.0} --- \textbf{tag} the image with a name and
  version
\item
  \texttt{.} --- the \textbf{build context} (current directory); Docker
  sends all files here to the build engine
\end{itemize}

List your images to confirm it was created:

\begin{Shaded}
\begin{Highlighting}[]
\NormalTok{docker image ls}
\end{Highlighting}
\end{Shaded}

Run a container from your custom image:

\begin{Shaded}
\begin{Highlighting}[]
\NormalTok{docker run {-}d {-}p 8080:80 {-}{-}name my{-}app my{-}nginx:1.0}
\end{Highlighting}
\end{Shaded}

Open \texttt{http://localhost:8080} --- you will see your custom page,
served directly from the image without any bind mounts.

\subsection{Understanding the Layer
Cache}\label{understanding-the-layer-cache}

Run \texttt{docker\ build} a second time without changing anything:

\begin{Shaded}
\begin{Highlighting}[]
\NormalTok{docker build {-}t my{-}nginx:1.0 .}
\end{Highlighting}
\end{Shaded}

You will see output like:

\begin{verbatim}
 => CACHED [1/2] FROM docker.io/library/nginx:alpine
 => CACHED [2/2] COPY html /usr/share/nginx/html
\end{verbatim}

Both steps are served from cache --- the build is instant. Now edit
\texttt{html/index.html} and rebuild:

\begin{Shaded}
\begin{Highlighting}[]
\NormalTok{docker build {-}t my{-}nginx:1.1 .}
\end{Highlighting}
\end{Shaded}

Only the \texttt{COPY} layer is rebuilt. The \texttt{FROM} layer
(downloading the base image) is still cached.

\textbf{Rule of thumb}: order your Dockerfile instructions from least to
most frequently changing. Put \texttt{RUN} commands that install
packages near the top, and \texttt{COPY} of your application code near
the bottom. This maximises cache reuse and keeps builds fast.

\subsection{Tagging and Versioning
Images}\label{tagging-and-versioning-images}

Good image tags tell you exactly what version is inside:

\begin{Shaded}
\begin{Highlighting}[]
\NormalTok{\# Tag with a version number}
\NormalTok{docker build {-}t my{-}nginx:1.0 .}

\NormalTok{\# Tag the same image as latest}
\NormalTok{docker tag my{-}nginx:1.0 my{-}nginx:latest}
\end{Highlighting}
\end{Shaded}

Always use explicit version tags in production. The \texttt{latest} tag
is convenient locally but can cause surprises in automated pipelines
because it changes silently.

\subsection{Cleaning Up}\label{cleaning-up}

Remove the running container:

\begin{Shaded}
\begin{Highlighting}[]
\NormalTok{docker rm {-}f my{-}app}
\end{Highlighting}
\end{Shaded}

Remove the image:

\begin{Shaded}
\begin{Highlighting}[]
\NormalTok{docker image rm my{-}nginx:1.0 my{-}nginx:1.1}
\end{Highlighting}
\end{Shaded}

Remove all unused images and build cache:

\begin{Shaded}
\begin{Highlighting}[]
\NormalTok{docker system prune {-}a}
\end{Highlighting}
\end{Shaded}

\subsection{Conclusion}\label{conclusion-2}

In this article we covered:

\begin{itemize}
\tightlist
\item
  The purpose of a \textbf{Dockerfile} and how it defines an image
\item
  The most important instructions: \texttt{FROM}, \texttt{COPY},
  \texttt{RUN}, \texttt{ENV}, \texttt{EXPOSE}, \texttt{CMD},
  \texttt{ENTRYPOINT}, \texttt{USER}
\item
  How to \textbf{build}, \textbf{tag}, and \textbf{run} a custom Nginx
  image
\item
  How the \textbf{layer cache} works and how to structure a Dockerfile
  to maximise it
\end{itemize}

The
\href{/Users/sasadangelo/github.com/sasadangelo/code4projects/how-docker-networking-works-2025/}{next
article} introduces Docker networking: how to connect containers
together so they can communicate over a private network.

\begin{center}\rule{0.5\linewidth}{0.5pt}\end{center}

If you enjoyed this article, don't forget to \textbf{give it a clap 👏},
\textbf{share it with your friends 🔗}, and \textbf{follow me for more
tips and tutorials on software development 📘}. Your support helps me
create more content like this --- thank you! 🙌

\section{How Docker Networking Works}\label{how-docker-networking-works}

\emph{Posted on ****}

\subsection{Introduction}\label{introduction-3}

This is the fourth article of the \textbf{Getting Started with Docker}
series. In the
\href{/Users/sasadangelo/github.com/sasadangelo/code4projects/dockerfile-and-building-custom-images/}{previous
article} we built a custom Nginx image. So far, each container has been
an isolated island. In a real application you often have multiple
containers that need to talk to each other --- a frontend container, a
backend container, a database.

This is where Docker networking comes in.

\subsection{Networking Overview in
Docker}\label{networking-overview-in-docker}

Docker has a pluggable networking system. The networking plugins are
called \textbf{drivers}. Docker provides these drivers out of the box:

\begin{enumerate}
\def\labelenumi{\arabic{enumi}.}
\tightlist
\item
  \textbf{bridge} --- the default; creates a private internal network on
  the host
\item
  \textbf{host} --- the container shares the host's network stack
  directly
\item
  \textbf{overlay} --- connects containers running on different hosts
  (used in Swarm mode)
\item
  \textbf{macvlan} --- assigns a real MAC address to a container, making
  it appear as a physical device on the network
\item
  \textbf{ipvlan} --- similar to macvlan but operates at layer 3; useful
  when the upstream switch does not allow multiple MACs per port
\item
  \textbf{none} --- networking is completely disabled
\end{enumerate}

For this series we focus on \textbf{bridge}, which is what you will use
for local development and single-host deployments.

\subsubsection{The bridge driver}\label{the-bridge-driver}

The \textbf{bridge} driver creates a private network on the host.
Containers attached to the same bridge can communicate with each other.
External access is granted by publishing ports with \texttt{-p}.

Behind the scenes, Docker creates a Linux bridge interface, assigns IP
addresses from a private subnet, and manages iptables rules to route
traffic in and out.

There are two kinds of bridge networks:

\begin{itemize}
\tightlist
\item
  \textbf{default bridge} (\texttt{bridge}) --- the one Docker uses
  automatically when you run a container without specifying a network.
  Containers can reach each other only by IP address, not by name.
\item
  \textbf{user-defined bridge} --- created with
  \texttt{docker\ network\ create}. Containers on the same user-defined
  bridge can \textbf{resolve each other by container name} via Docker's
  built-in DNS. This is the right choice for any multi-container setup.
\end{itemize}

\subsubsection{The host driver}\label{the-host-driver}

With the \textbf{host} driver the container skips Docker's network layer
entirely and shares the host's network interfaces. Port mapping
(\texttt{-p}) is not needed --- if the container listens on port 80, it
is directly reachable on the host's port 80.

Useful for performance-sensitive workloads or tools that need to monitor
host network traffic. Note: \texttt{-\/-network\ host} does not work on
Docker Desktop for Mac or Windows --- it only works natively on Linux.

\subsubsection{The overlay driver}\label{the-overlay-driver}

The \textbf{overlay} driver enables containers on different physical
hosts to communicate as if they were on the same network. It is the
networking backbone for \textbf{Docker Swarm} clusters. We will cover
this when we discuss Swarm in a future article.

\subsubsection{The macvlan driver}\label{the-macvlan-driver}

The \textbf{macvlan} driver assigns a MAC address to a container. This
makes it appear as a physical device directly on the network segment. It
is the right choice for legacy applications that expect to be directly
connected to the physical network rather than routed through Docker's
network stack.

\subsection{Docker Networking
Commands}\label{docker-networking-commands}

Create a user-defined bridge network:

\begin{Shaded}
\begin{Highlighting}[]
\NormalTok{docker network create my{-}network}
\end{Highlighting}
\end{Shaded}

To specify a subnet and gateway explicitly:

\begin{Shaded}
\begin{Highlighting}[]
\NormalTok{docker network create \textbackslash{}}
\NormalTok{  {-}d bridge \textbackslash{}}
\NormalTok{  {-}{-}subnet=10.10.1.0/24 \textbackslash{}}
\NormalTok{  {-}{-}gateway=10.10.1.1 \textbackslash{}}
\NormalTok{  my{-}network}
\end{Highlighting}
\end{Shaded}

List all networks:

\begin{Shaded}
\begin{Highlighting}[]
\NormalTok{docker network ls}
\end{Highlighting}
\end{Shaded}

Inspect a network (shows connected containers, subnet, gateway):

\begin{Shaded}
\begin{Highlighting}[]
\NormalTok{docker network inspect my{-}network}
\end{Highlighting}
\end{Shaded}

Remove a network:

\begin{Shaded}
\begin{Highlighting}[]
\NormalTok{docker network rm my{-}network}
\end{Highlighting}
\end{Shaded}

Connect a running container to a network:

\begin{Shaded}
\begin{Highlighting}[]
\NormalTok{docker network connect my{-}network my{-}container}
\end{Highlighting}
\end{Shaded}

\subsection{Practical Example: Frontend and Backend on the Same
Network}\label{practical-example-frontend-and-backend-on-the-same-network}

Let us build a concrete example. We will run two containers on the same
user-defined bridge network and have them communicate by name --- no IP
addresses needed.

\subsubsection{Create the network}\label{create-the-network}

\begin{Shaded}
\begin{Highlighting}[]
\NormalTok{docker network create app{-}network}
\end{Highlighting}
\end{Shaded}

\subsubsection{Start the ``backend''
container}\label{start-the-backend-container}

For simplicity we use another Nginx instance as a mock backend serving a
JSON response.

\begin{Shaded}
\begin{Highlighting}[]
\NormalTok{mkdir backend}
\NormalTok{cat \textgreater{} backend/index.html  frontend/index.html }

\NormalTok{  Frontend}
  
\NormalTok{    Frontend container}
\NormalTok{    Backend is reachable at http://backend}
  

\NormalTok{EOF}

\NormalTok{docker run {-}d \textbackslash{}}
\NormalTok{  {-}{-}name frontend \textbackslash{}}
\NormalTok{  {-}{-}network app{-}network \textbackslash{}}
\NormalTok{  {-}p 8080:80 \textbackslash{}}
\NormalTok{  {-}v $(pwd)/frontend:/usr/share/nginx/html:ro \textbackslash{}}
\NormalTok{  nginx:alpine}
\end{Highlighting}
\end{Shaded}

\subsubsection{Verify containers can resolve each other by
name}\label{verify-containers-can-resolve-each-other-by-name}

\begin{Shaded}
\begin{Highlighting}[]
\NormalTok{docker exec frontend ping {-}c 3 backend}
\end{Highlighting}
\end{Shaded}

Because both containers are on the same user-defined bridge network,
Docker's built-in DNS resolves \texttt{backend} to the correct container
IP automatically. No need to hardcode IP addresses anywhere.

Open \texttt{http://localhost:8080} in your browser to see the frontend.
The backend is reachable internally at \texttt{http://backend} from
within the \texttt{frontend} container.

\subsubsection{Clean up}\label{clean-up}

\begin{Shaded}
\begin{Highlighting}[]
\NormalTok{docker rm {-}f frontend backend}
\NormalTok{docker network rm app{-}network}
\end{Highlighting}
\end{Shaded}

\subsection{Conclusion}\label{conclusion-3}

In this article we covered:

\begin{itemize}
\tightlist
\item
  Docker's five network drivers: \textbf{bridge}, \textbf{host},
  \textbf{overlay}, \textbf{macvlan}, \textbf{ipvlan}, and \textbf{none}
\item
  The difference between the \textbf{default bridge} and
  \textbf{user-defined bridge} networks
\item
  How \textbf{automatic DNS resolution} works on user-defined bridges
  (containers talk by name, not IP)
\item
  Essential networking commands: \texttt{create}, \texttt{ls},
  \texttt{inspect}, \texttt{rm}, \texttt{connect}
\item
  A practical example with frontend and backend containers communicating
  by name
\end{itemize}

The
\href{/Users/sasadangelo/github.com/sasadangelo/code4projects/how-docker-volumes-work/}{next
article} covers \textbf{Docker volumes}: how to persist data outside a
container so it survives restarts and replacements.

\begin{center}\rule{0.5\linewidth}{0.5pt}\end{center}

If you enjoyed this article, don't forget to \textbf{give it a clap 👏},
\textbf{share it with your friends 🔗}, and \textbf{follow me for more
tips and tutorials on software development 📘}. Your support helps me
create more content like this --- thank you! 🙌

\section{How Docker Volumes Work}\label{how-docker-volumes-work}

\emph{Posted on ****}

\subsection{Introduction}\label{introduction-4}

This is the fifth article of the \textbf{Getting Started with Docker}
series. In the
\href{/Users/sasadangelo/github.com/sasadangelo/code4projects/how-docker-networking-works-2025/}{previous
article} we connected containers over a private network. Now we need to
address another fundamental question: \textbf{what happens to your data
when a container is stopped or replaced?}

By default, a container's filesystem is \textbf{ephemeral} ---
everything written inside it disappears when the container is removed.
This is fine for stateless services like a web server, but it is a
problem as soon as you have anything that needs to survive: logs,
configuration, user data, database files.

Docker solves this with \textbf{volumes}.

\subsection{The Upgrade Problem}\label{the-upgrade-problem}

One of Docker's main benefits is that upgrading an application is as
simple as replacing the old container with a new one. But if your data
is stored inside the container, replacing it means losing the data.

The solution is to \textbf{separate the application binaries from the
data}. Store the binaries in the image and the data in a volume that
lives outside the container. When you replace the container, the volume
stays where it is and the new container picks it up.

\subsection{Docker Storage Types}\label{docker-storage-types}

Docker provides three options for persisting data from a container to
the host filesystem.

\subsubsection{Volumes}\label{volumes}

\textbf{Volumes} are managed entirely by Docker. They are stored in a
directory on the host that Docker controls (typically
\texttt{/var/lib/docker/volumes/} on Linux) and you interact with them
through Docker CLI commands --- not by poking around the filesystem
directly.

Volumes are the \textbf{recommended way to persist data} in Docker. From
the \href{https://docs.docker.com/storage/volumes/}{official
documentation}, some of the reasons:

\begin{itemize}
\tightlist
\item
  Easier to back up and migrate than bind mounts
\item
  Manageable through the Docker CLI and API
\item
  Work on both Linux and Windows containers
\item
  Can be safely shared among multiple containers
\item
  Support remote storage drivers (NFS, cloud block storage, etc.)
\item
  New volumes can be pre-populated with data from a container
\end{itemize}

\subsubsection{Bind Mounts}\label{bind-mounts}

\textbf{Bind mounts} map a specific directory or file on the
\textbf{host filesystem} into the container. Unlike volumes, bind mounts
are not managed by Docker --- any process on the host can read and write
them.

Bind mounts are ideal when:

\begin{itemize}
\tightlist
\item
  You want to share source code between your editor on the host and the
  running container (live reload during development)
\item
  You need the storage to be on a specific filesystem type
\item
  You need direct access to the files from the host without going
  through Docker
\end{itemize}

We already used a bind mount in the first article:
\texttt{-v\ \$(pwd)/index.html:/usr/share/nginx/html/index.html:ro}.

\subsubsection{Tmpfs Mounts}\label{tmpfs-mounts}

\textbf{Tmpfs mounts} store data in the host's memory, not on disk. The
data is never written to disk and disappears when the container stops.
They are useful for sensitive data (secrets, tokens) that should never
touch disk, or for temporary scratch space that requires very fast I/O.

Tmpfs mounts are only available on Linux hosts.

\subsection{\texorpdfstring{The \texttt{-\/-mount}
Syntax}{The -\/-mount Syntax}}\label{the---mount-syntax}

Docker has two ways to attach storage to a container: the older
\texttt{-v} flag and the newer \texttt{-\/-mount} flag. The
\texttt{-\/-mount} syntax is more explicit and is the one recommended by
the official documentation today.

Comparison:

\begin{Shaded}
\begin{Highlighting}[]
\NormalTok{\# Old syntax ({-}v)}
\NormalTok{docker run {-}v my{-}volume:/app/data nginx:alpine}

\NormalTok{\# New syntax ({-}{-}mount) — equivalent}
\NormalTok{docker run {-}{-}mount type=volume,source=my{-}volume,target=/app/data nginx:alpine}
\end{Highlighting}
\end{Shaded}

The \texttt{-\/-mount} form accepts key-value pairs:

\begin{itemize}
\tightlist
\item
  \texttt{type} --- \texttt{volume}, \texttt{bind}, or \texttt{tmpfs}
\item
  \texttt{source} (or \texttt{src}) --- volume name or host path (for
  bind mounts)
\item
  \texttt{target} (or \texttt{dst}) --- path inside the container
\item
  \texttt{readonly} --- makes the mount read-only
\end{itemize}

Both syntaxes work; \texttt{-\/-mount} is preferred in scripts and
documentation for clarity.

\subsection{Docker Volumes Cheat
Sheet}\label{docker-volumes-cheat-sheet}

Create a named volume:

\begin{Shaded}
\begin{Highlighting}[]
\NormalTok{docker volume create my{-}data}
\end{Highlighting}
\end{Shaded}

List all volumes:

\begin{Shaded}
\begin{Highlighting}[]
\NormalTok{docker volume ls}
\end{Highlighting}
\end{Shaded}

Inspect a volume (shows the mount point on the host):

\begin{Shaded}
\begin{Highlighting}[]
\NormalTok{docker volume inspect my{-}data}
\end{Highlighting}
\end{Shaded}

Remove a volume:

\begin{Shaded}
\begin{Highlighting}[]
\NormalTok{docker volume rm my{-}data}
\end{Highlighting}
\end{Shaded}

Remove all unused volumes:

\begin{Shaded}
\begin{Highlighting}[]
\NormalTok{docker volume prune}
\end{Highlighting}
\end{Shaded}

\subsection{Practical Example: Nginx with Persistent
Content}\label{practical-example-nginx-with-persistent-content}

Let us demonstrate volumes with our Nginx container. We will create a
volume, write content into it from a temporary container, then serve
that content with Nginx --- and show that the content survives a
container replacement.

\subsubsection{Create the volume}\label{create-the-volume}

\begin{Shaded}
\begin{Highlighting}[]
\NormalTok{docker volume create web{-}content}
\end{Highlighting}
\end{Shaded}

\subsubsection{Populate the volume}\label{populate-the-volume}

\begin{Shaded}
\begin{Highlighting}[]
\NormalTok{docker run {-}{-}rm \textbackslash{}}
\NormalTok{  {-}{-}mount type=volume,source=web{-}content,target=/data \textbackslash{}}
\NormalTok{  alpine sh {-}c \textquotesingle{}echo "Content from a volume" \textgreater{} /data/index.html\textquotesingle{}}
\end{Highlighting}
\end{Shaded}

This starts a temporary Alpine container, mounts the volume at
\texttt{/data}, writes a file, and immediately exits and removes itself
(\texttt{-\/-rm}). The volume and its content persist.

\subsubsection{Serve the content with
Nginx}\label{serve-the-content-with-nginx}

\begin{Shaded}
\begin{Highlighting}[]
\NormalTok{docker run {-}d \textbackslash{}}
\NormalTok{  {-}{-}name nginx{-}volume \textbackslash{}}
\NormalTok{  {-}p 8080:80 \textbackslash{}}
\NormalTok{  {-}{-}mount type=volume,source=web{-}content,target=/usr/share/nginx/html \textbackslash{}}
\NormalTok{  nginx:alpine}
\end{Highlighting}
\end{Shaded}

Open \texttt{http://localhost:8080} --- you see the page from the
volume.

\subsubsection{Replace the container --- data
survives}\label{replace-the-container-data-survives}

\begin{Shaded}
\begin{Highlighting}[]
\NormalTok{docker rm {-}f nginx{-}volume}

\NormalTok{docker run {-}d \textbackslash{}}
\NormalTok{  {-}{-}name nginx{-}volume{-}new \textbackslash{}}
\NormalTok{  {-}p 8080:80 \textbackslash{}}
\NormalTok{  {-}{-}mount type=volume,source=web{-}content,target=/usr/share/nginx/html \textbackslash{}}
\NormalTok{  nginx:alpine}
\end{Highlighting}
\end{Shaded}

Open \texttt{http://localhost:8080} again. The same page is still there.
The volume outlived the container.

\subsubsection{Clean up}\label{clean-up-1}

\begin{Shaded}
\begin{Highlighting}[]
\NormalTok{docker rm {-}f nginx{-}volume{-}new}
\NormalTok{docker volume rm web{-}content}
\end{Highlighting}
\end{Shaded}

\subsection{Volumes vs Bind Mounts: When to Use
Which}\label{volumes-vs-bind-mounts-when-to-use-which}

\begin{longtable}[]{@{}
  >{\raggedright\arraybackslash}p{(\linewidth - 4\tabcolsep) * \real{0.3333}}
  >{\raggedright\arraybackslash}p{(\linewidth - 4\tabcolsep) * \real{0.3333}}
  >{\raggedright\arraybackslash}p{(\linewidth - 4\tabcolsep) * \real{0.3333}}@{}}
\toprule\noalign{}
\begin{minipage}[b]{\linewidth}\raggedright
\end{minipage} & \begin{minipage}[b]{\linewidth}\raggedright
Volumes
\end{minipage} & \begin{minipage}[b]{\linewidth}\raggedright
Bind Mounts
\end{minipage} \\
\midrule\noalign{}
\endhead
\bottomrule\noalign{}
\endlastfoot
Managed by Docker & Yes & No \\
Portable across machines & Yes (with volume drivers) & No (path must
exist on host) \\
Best for production data & Yes & No \\
Best for dev code sharing & No & Yes \\
Access from host directly & Via \texttt{docker\ volume\ inspect} &
Direct filesystem path \\
Works on Windows containers & Yes & Limited \\
\end{longtable}

\subsection{Conclusion}\label{conclusion-4}

In this article we covered:

\begin{itemize}
\tightlist
\item
  Why containers are ephemeral by default and why that matters
\item
  The three Docker storage types: \textbf{volumes}, \textbf{bind
  mounts}, and \textbf{tmpfs}
\item
  The modern \textbf{\texttt{-\/-mount}} syntax vs the older \texttt{-v}
  flag
\item
  Essential volume commands: \texttt{create}, \texttt{ls},
  \texttt{inspect}, \texttt{rm}, \texttt{prune}
\item
  A practical example showing that volume data survives container
  replacement
\item
  A comparison table to decide when to use volumes vs bind mounts
\end{itemize}

The
\href{/Users/sasadangelo/github.com/sasadangelo/code4projects/how-docker-compose-works-2025/}{next
article} introduces \textbf{Docker Compose}: how to manage a
multi-container application with a single YAML file instead of juggling
multiple \texttt{docker\ run} commands.

\begin{center}\rule{0.5\linewidth}{0.5pt}\end{center}

If you enjoyed this article, don't forget to \textbf{give it a clap 👏},
\textbf{share it with your friends 🔗}, and \textbf{follow me for more
tips and tutorials on software development 📘}. Your support helps me
create more content like this --- thank you! 🙌

\section{How Docker Compose Works}\label{how-docker-compose-works}

\emph{Posted on ****}

\subsection{Introduction}\label{introduction-5}

This is the sixth article of the \textbf{Getting Started with Docker}
series. In the
\href{/Users/sasadangelo/github.com/sasadangelo/code4projects/how-docker-volumes-work/}{previous
article} we learned how to persist data with volumes. So far, to run our
application we have been juggling multiple \texttt{docker\ run} commands
--- one for the network, one for each container, one for each volume.
This quickly becomes hard to manage and impossible to share reliably
with a team.

\textbf{Docker Compose} solves this: you describe your entire
application in a single YAML file and launch everything with one
command.

\subsection{What is Docker Compose?}\label{what-is-docker-compose}

Docker Compose is a tool for defining and running multi-container Docker
applications. You describe your services, networks, and volumes in a
\texttt{docker-compose.yml} file, and Docker Compose translates that
into the correct sequence of \texttt{docker} commands.

\begin{quote}
\textbf{Important}: the original \texttt{docker-compose} command (v1,
written in Python) was deprecated in May 2023 and is no longer
maintained. The current tool is \textbf{\texttt{docker\ compose}} (v2,
written in Go), which is a plugin built directly into the Docker CLI.
Always use \texttt{docker\ compose} (with a space, not a hyphen). If you
are on Podman, use \texttt{podman\ compose} which is compatible with the
same YAML format.
\end{quote}

You can think of \texttt{docker\ compose} as a wrapper that orchestrates
all the \texttt{docker} commands you would otherwise run manually ---
but declaratively, from a single file.

\subsection{The docker-compose.yml
File}\label{the-docker-compose.yml-file}

The Compose file is a YAML document. Its top-level keys are:

\begin{itemize}
\tightlist
\item
  \texttt{services} --- the containers that make up your application
\item
  \texttt{volumes} --- named volumes used by the services
\item
  \texttt{networks} --- custom networks (if not specified, Compose
  creates a default bridge network automatically)
\end{itemize}

\subsubsection{\texorpdfstring{A Note on the \texttt{version}
Field}{A Note on the version Field}}\label{a-note-on-the-version-field}

Older tutorials and documentation show a
\texttt{version:\ \textquotesingle{}3.x\textquotesingle{}} field at the
top of the Compose file. This field is \textbf{deprecated} and should be
omitted. Modern Docker Compose uses the
\href{https://compose-spec.io/}{Compose Specification} which does not
require a version declaration.

\subsection{Practical Example: Frontend and Backend with
Compose}\label{practical-example-frontend-and-backend-with-compose}

We will rebuild the two-container application from the networking
article using Docker Compose, adding a named volume for the frontend
content.

\subsubsection{Project structure}\label{project-structure}

\begin{verbatim}
my-app/
├── docker-compose.yml
├── frontend/
│   └── index.html
└── backend/
    └── index.html
\end{verbatim}

Create the directories and content:

\begin{Shaded}
\begin{Highlighting}[]
\NormalTok{mkdir {-}p my{-}app/frontend my{-}app/backend}
\NormalTok{cd my{-}app}

\NormalTok{cat \textgreater{} frontend/index.html }

\NormalTok{  Frontend}
  
\NormalTok{    Frontend}
\NormalTok{    This app is managed by Docker Compose.}
  

\NormalTok{EOF}

\NormalTok{cat \textgreater{} backend/index.html \textless{}\textless{} \textquotesingle{}EOF\textquotesingle{}}
\NormalTok{\{"status": "ok", "message": "Hello from the backend!"\}}
\NormalTok{EOF}
\end{Highlighting}
\end{Shaded}

\subsubsection{The docker-compose.yml
file}\label{the-docker-compose.yml-file-1}

\begin{Shaded}
\begin{Highlighting}[]
\FunctionTok{services}\KeywordTok{:}
\AttributeTok{  }\FunctionTok{frontend}\KeywordTok{:}
\AttributeTok{    }\FunctionTok{image}\KeywordTok{:}\AttributeTok{ nginx:alpine}
\AttributeTok{    }\FunctionTok{container\_name}\KeywordTok{:}\AttributeTok{ frontend}
\AttributeTok{    }\FunctionTok{ports}\KeywordTok{:}
\AttributeTok{      }\KeywordTok{{-}}\AttributeTok{ }\StringTok{"8080:80"}
\AttributeTok{    }\FunctionTok{volumes}\KeywordTok{:}
\AttributeTok{      }\KeywordTok{{-}}\AttributeTok{ ./frontend:/usr/share/nginx/html:ro}
\AttributeTok{    }\FunctionTok{networks}\KeywordTok{:}
\AttributeTok{      }\KeywordTok{{-}}\AttributeTok{ app{-}network}
\AttributeTok{    }\FunctionTok{depends\_on}\KeywordTok{:}
\AttributeTok{      }\KeywordTok{{-}}\AttributeTok{ backend}

\AttributeTok{  }\FunctionTok{backend}\KeywordTok{:}
\AttributeTok{    }\FunctionTok{image}\KeywordTok{:}\AttributeTok{ nginx:alpine}
\AttributeTok{    }\FunctionTok{container\_name}\KeywordTok{:}\AttributeTok{ backend}
\AttributeTok{    }\FunctionTok{volumes}\KeywordTok{:}
\AttributeTok{      }\KeywordTok{{-}}\AttributeTok{ ./backend:/usr/share/nginx/html:ro}
\AttributeTok{    }\FunctionTok{networks}\KeywordTok{:}
\AttributeTok{      }\KeywordTok{{-}}\AttributeTok{ app{-}network}

\FunctionTok{networks}\KeywordTok{:}
\AttributeTok{  }\FunctionTok{app{-}network}\KeywordTok{:}
\AttributeTok{    }\FunctionTok{driver}\KeywordTok{:}\AttributeTok{ bridge}
\end{Highlighting}
\end{Shaded}

Key things to notice:

\begin{itemize}
\tightlist
\item
  \textbf{No \texttt{version:} field} --- not needed with modern Compose
\item
  \textbf{\texttt{depends\_on}} --- ensures the backend container starts
  before the frontend. For health-based ordering, you can use
  \texttt{condition:\ service\_healthy} together with a
  \texttt{healthcheck} definition
\item
  \textbf{Networks are created automatically} --- Compose creates
  \texttt{app-network} for you
\item
  \textbf{Bind mounts} use relative paths from the directory containing
  \texttt{docker-compose.yml}
\end{itemize}

\subsection{Docker Compose Commands}\label{docker-compose-commands}

The most important commands:

\begin{Shaded}
\begin{Highlighting}[]
\NormalTok{\# Start all services (build images if needed, create networks and volumes)}
\NormalTok{docker compose up {-}d}

\NormalTok{\# Stop and remove containers, networks (volumes are preserved)}
\NormalTok{docker compose down}

\NormalTok{\# Stop and remove everything including volumes}
\NormalTok{docker compose down {-}v}

\NormalTok{\# View logs of all services}
\NormalTok{docker compose logs}

\NormalTok{\# Follow logs in real time}
\NormalTok{docker compose logs {-}f}

\NormalTok{\# List running containers managed by this Compose file}
\NormalTok{docker compose ps}

\NormalTok{\# Execute a command inside a running service container}
\NormalTok{docker compose exec frontend sh}

\NormalTok{\# Rebuild images (useful after changing a Dockerfile)}
\NormalTok{docker compose build}

\NormalTok{\# Pull the latest versions of all images}
\NormalTok{docker compose pull}

\NormalTok{\# Restart a single service}
\NormalTok{docker compose restart frontend}
\end{Highlighting}
\end{Shaded}

The key difference from \texttt{docker} commands is that all of these
act on the \textbf{entire application} defined in the Compose file, not
a single container.

\subsection{Environment Variables}\label{environment-variables}

Hard-coded values in \texttt{docker-compose.yml} are fine for
development, but for anything sensitive (passwords, API keys) or
environment-specific (URLs, port numbers) you should use variables.

Docker Compose automatically reads a file named \texttt{.env} in the
same directory:

\begin{Shaded}
\begin{Highlighting}[]
\NormalTok{\# .env}
\NormalTok{NGINX\_PORT=8080}
\NormalTok{APP\_ENV=development}
\end{Highlighting}
\end{Shaded}

Then reference them in \texttt{docker-compose.yml}:

\begin{Shaded}
\begin{Highlighting}[]
\FunctionTok{services}\KeywordTok{:}
\AttributeTok{  }\FunctionTok{frontend}\KeywordTok{:}
\AttributeTok{    }\FunctionTok{ports}\KeywordTok{:}
\AttributeTok{      }\KeywordTok{{-}}\AttributeTok{ }\StringTok{"$\{NGINX\_PORT\}:80"}
\AttributeTok{    }\FunctionTok{environment}\KeywordTok{:}
\AttributeTok{      }\KeywordTok{{-}}\AttributeTok{ APP\_ENV=$\{APP\_ENV\}}
\end{Highlighting}
\end{Shaded}

Never commit \texttt{.env} files containing secrets to version control.
Add \texttt{.env} to your \texttt{.gitignore}.

\subsection{Starting the Application}\label{starting-the-application}

From the \texttt{my-app/} directory:

\begin{Shaded}
\begin{Highlighting}[]
\NormalTok{docker compose up {-}d}
\end{Highlighting}
\end{Shaded}

Compose will:

\begin{enumerate}
\def\labelenumi{\arabic{enumi}.}
\tightlist
\item
  Create the \texttt{app-network} bridge network
\item
  Start the \texttt{backend} container
\item
  Start the \texttt{frontend} container
\item
  Map port 8080 on your host to port 80 on the frontend container
\end{enumerate}

Open \texttt{http://localhost:8080} --- the frontend is running. Check
that both services are up:

\begin{Shaded}
\begin{Highlighting}[]
\NormalTok{docker compose ps}
\end{Highlighting}
\end{Shaded}

Verify the backend is reachable by name from the frontend:

\begin{Shaded}
\begin{Highlighting}[]
\NormalTok{docker compose exec frontend ping {-}c 3 backend}
\end{Highlighting}
\end{Shaded}

\subsubsection{Clean up}\label{clean-up-2}

\begin{Shaded}
\begin{Highlighting}[]
\NormalTok{docker compose down}
\end{Highlighting}
\end{Shaded}

\subsection{Conclusion}\label{conclusion-5}

In this article we covered:

\begin{itemize}
\tightlist
\item
  What Docker Compose is and why it replaces manual \texttt{docker\ run}
  scripts
\item
  The difference between \texttt{docker-compose} v1 (deprecated) and
  \textbf{\texttt{docker\ compose}} v2 (current)
\item
  The structure of a \texttt{docker-compose.yml} file:
  \texttt{services}, \texttt{volumes}, \texttt{networks}
\item
  Why the \texttt{version:} field is no longer needed
\item
  Essential Compose commands: \texttt{up}, \texttt{down}, \texttt{logs},
  \texttt{ps}, \texttt{exec}, \texttt{build}
\item
  How to manage configuration with \textbf{environment variables} and
  \texttt{.env} files
\end{itemize}

The
\href{/Users/sasadangelo/github.com/sasadangelo/code4projects/docker-security-best-practices/}{next
article} closes the series with \textbf{Docker security best practices}:
how to run containers safely, limit their privileges, and scan images
for vulnerabilities.

\begin{center}\rule{0.5\linewidth}{0.5pt}\end{center}

If you enjoyed this article, don't forget to \textbf{give it a clap 👏},
\textbf{share it with your friends 🔗}, and \textbf{follow me for more
tips and tutorials on software development 📘}. Your support helps me
create more content like this --- thank you! 🙌

\section{Docker Security Best
Practices}\label{docker-security-best-practices}

\emph{Posted on ****}

\subsection{Introduction}\label{introduction-6}

This is the seventh and final article of the \textbf{Getting Started
with Docker} series. In the
\href{/Users/sasadangelo/github.com/sasadangelo/code4projects/how-docker-compose-works-2025/}{previous
article} we orchestrated a multi-container application with Docker
Compose. Now we turn our attention to \textbf{security}.

Containers are isolated --- but isolation is not the same as security. A
container running as root, with a writable filesystem and full Linux
capabilities, is a significant risk if it is ever compromised. The good
news is that Docker provides several mechanisms to reduce the attack
surface, and most of them require only a few lines of configuration.

\subsection{1. Never Run as Root}\label{never-run-as-root}

By default, processes inside a container run as \textbf{root} (UID 0).
This means that if an attacker exploits a vulnerability in your
application, they have root access inside the container --- and
depending on how the container is configured, they may be able to escape
to the host.

The fix is simple: add a \texttt{USER} instruction to your Dockerfile.

\begin{Shaded}
\begin{Highlighting}[]
\KeywordTok{FROM}\NormalTok{ nginx:alpine}

\KeywordTok{COPY}\NormalTok{ html /usr/share/nginx/html}

\CommentTok{\# Switch to the non{-}root user that the nginx image already provides}
\KeywordTok{USER}\NormalTok{ nginx}

\KeywordTok{EXPOSE}\NormalTok{ 80}
\KeywordTok{CMD}\NormalTok{ [}\StringTok{"nginx"}\NormalTok{, }\StringTok{"{-}g"}\NormalTok{, }\StringTok{"daemon off;"}\NormalTok{]}
\end{Highlighting}
\end{Shaded}

The official \texttt{nginx:alpine} image already includes a non-root
\texttt{nginx} user. For your own images, create a dedicated user:

\begin{Shaded}
\begin{Highlighting}[]
\KeywordTok{FROM}\NormalTok{ alpine:3.20}

\KeywordTok{RUN} \ExtensionTok{addgroup} \AttributeTok{{-}S}\NormalTok{ appgroup }\KeywordTok{\&\&} \ExtensionTok{adduser} \AttributeTok{{-}S}\NormalTok{ appuser }\AttributeTok{{-}G}\NormalTok{ appgroup}

\KeywordTok{COPY} \OperatorTok{{-}{-}chown=appuser:appgroup}\NormalTok{ app /app}

\KeywordTok{USER}\NormalTok{ appuser}

\KeywordTok{CMD}\NormalTok{ [}\StringTok{"/app/start.sh"}\NormalTok{]}
\end{Highlighting}
\end{Shaded}

The \texttt{-\/-chown} flag in \texttt{COPY} ensures the files are owned
by the correct user from the start.

You can also enforce non-root at runtime without changing the
Dockerfile:

\begin{Shaded}
\begin{Highlighting}[]
\NormalTok{docker run {-}{-}user 1001:1001 nginx:alpine}
\end{Highlighting}
\end{Shaded}

\subsection{2. Use a Read-Only
Filesystem}\label{use-a-read-only-filesystem}

If your application does not need to write to its own filesystem, make
it read-only:

\begin{Shaded}
\begin{Highlighting}[]
\NormalTok{docker run {-}{-}read{-}only {-}p 8080:80 nginx:alpine}
\end{Highlighting}
\end{Shaded}

If the application needs to write to specific paths (temp files, logs),
mount those paths as \texttt{tmpfs}:

\begin{Shaded}
\begin{Highlighting}[]
\NormalTok{docker run {-}{-}read{-}only \textbackslash{}}
\NormalTok{  {-}{-}tmpfs /tmp \textbackslash{}}
\NormalTok{  {-}{-}tmpfs /var/cache/nginx \textbackslash{}}
\NormalTok{  {-}{-}tmpfs /var/run \textbackslash{}}
\NormalTok{  {-}p 8080:80 \textbackslash{}}
\NormalTok{  nginx:alpine}
\end{Highlighting}
\end{Shaded}

A read-only filesystem means that even if an attacker gains code
execution inside the container, they cannot modify the application
binaries or install tools.

In Docker Compose:

\begin{Shaded}
\begin{Highlighting}[]
\FunctionTok{services}\KeywordTok{:}
\AttributeTok{  }\FunctionTok{frontend}\KeywordTok{:}
\AttributeTok{    }\FunctionTok{image}\KeywordTok{:}\AttributeTok{ nginx:alpine}
\AttributeTok{    }\FunctionTok{read\_only}\KeywordTok{:}\AttributeTok{ }\CharTok{true}
\AttributeTok{    }\FunctionTok{tmpfs}\KeywordTok{:}
\AttributeTok{      }\KeywordTok{{-}}\AttributeTok{ /tmp}
\AttributeTok{      }\KeywordTok{{-}}\AttributeTok{ /var/cache/nginx}
\AttributeTok{      }\KeywordTok{{-}}\AttributeTok{ /var/run}
\end{Highlighting}
\end{Shaded}

\subsection{3. Drop Linux Capabilities}\label{drop-linux-capabilities}

Linux capabilities are fine-grained divisions of root privileges. By
default, Docker grants a container a subset of them. You can drop all
capabilities and add back only what your application actually needs:

\begin{Shaded}
\begin{Highlighting}[]
\NormalTok{docker run {-}{-}cap{-}drop=ALL {-}{-}cap{-}add=NET\_BIND\_SERVICE nginx:alpine}
\end{Highlighting}
\end{Shaded}

\texttt{NET\_BIND\_SERVICE} allows binding to ports below 1024 (like
port 80). Most application containers do not need any other capability.

In Docker Compose:

\begin{Shaded}
\begin{Highlighting}[]
\FunctionTok{services}\KeywordTok{:}
\AttributeTok{  }\FunctionTok{frontend}\KeywordTok{:}
\AttributeTok{    }\FunctionTok{image}\KeywordTok{:}\AttributeTok{ nginx:alpine}
\AttributeTok{    }\FunctionTok{cap\_drop}\KeywordTok{:}
\AttributeTok{      }\KeywordTok{{-}}\AttributeTok{ ALL}
\AttributeTok{    }\FunctionTok{cap\_add}\KeywordTok{:}
\AttributeTok{      }\KeywordTok{{-}}\AttributeTok{ NET\_BIND\_SERVICE}
\end{Highlighting}
\end{Shaded}

\subsection{4. Never Put Secrets in Environment Variables (or
Images)}\label{never-put-secrets-in-environment-variables-or-images}

A common mistake is passing passwords or API keys through environment
variables in \texttt{docker-compose.yml} or baking them into the image
with \texttt{ENV}. Both approaches expose secrets in
\texttt{docker\ inspect} output and in image layers.

\textbf{What to do instead}:

\begin{itemize}
\tightlist
\item
  Use \textbf{Docker secrets} (available in Swarm mode) which mount
  secrets as files in \texttt{/run/secrets/}
\item
  Use a \textbf{secrets manager} (HashiCorp Vault, AWS Secrets Manager,
  etc.)
\item
  As a minimum, use a \texttt{.env} file that is \textbf{never committed
  to version control} and reference it in Compose
\end{itemize}

\begin{Shaded}
\begin{Highlighting}[]
\CommentTok{\# docker{-}compose.yml}
\FunctionTok{services}\KeywordTok{:}
\AttributeTok{  }\FunctionTok{app}\KeywordTok{:}
\AttributeTok{    }\FunctionTok{image}\KeywordTok{:}\AttributeTok{ my{-}app}
\AttributeTok{    }\FunctionTok{env\_file}\KeywordTok{:}
\AttributeTok{      }\KeywordTok{{-}}\AttributeTok{ .env.secrets}\CommentTok{   \# never committed to git}
\end{Highlighting}
\end{Shaded}

Add to \texttt{.gitignore}:

\begin{verbatim}
.env
.env.secrets
*.env
\end{verbatim}

\subsection{5. Use Minimal Base Images}\label{use-minimal-base-images}

Every package in your base image is a potential attack surface. Prefer
minimal base images:

\begin{itemize}
\tightlist
\item
  \texttt{alpine} --- \textasciitilde5MB, very popular, uses musl libc
\item
  \texttt{distroless} (Google) --- no shell, no package manager, only
  the runtime and your app
\item
  \texttt{scratch} --- completely empty, for statically compiled
  binaries (Go, Rust)
\end{itemize}

For our Nginx example we already use \texttt{nginx:alpine}. Going even
smaller would mean building a custom Nginx binary, which is rarely
necessary.

The rule: \textbf{if it is not needed, it should not be in the image}.

\subsection{6. Scan Images for Known
Vulnerabilities}\label{scan-images-for-known-vulnerabilities}

Even minimal images can contain packages with known CVEs. Docker
provides a built-in scanner called \textbf{Docker Scout}:

\begin{Shaded}
\begin{Highlighting}[]
\NormalTok{\# Quick vulnerability overview}
\NormalTok{docker scout quickview nginx:alpine}

\NormalTok{\# Full CVE list}
\NormalTok{docker scout cves nginx:alpine}

\NormalTok{\# Recommendations to fix vulnerabilities}
\NormalTok{docker scout recommendations nginx:alpine}
\end{Highlighting}
\end{Shaded}

Docker Scout is available in Docker Desktop and as a CLI plugin. It
integrates with Docker Hub and can be added to CI/CD pipelines to fail a
build when high-severity vulnerabilities are detected.

An alternative open-source scanner is \textbf{Trivy}:

\begin{Shaded}
\begin{Highlighting}[]
\NormalTok{trivy image nginx:alpine}
\end{Highlighting}
\end{Shaded}

Make image scanning part of your build pipeline, not an afterthought.

\subsection{7. Keep Images Up to Date}\label{keep-images-up-to-date}

A scanned image with no vulnerabilities today may have CVEs tomorrow.
Pin your base image to a specific digest rather than a floating tag, and
automate rebuild checks:

\begin{Shaded}
\begin{Highlighting}[]
\CommentTok{\# Pin to a specific digest (reproducible, auditable)}
\KeywordTok{FROM}\NormalTok{ nginx:alpine@sha256:...}
\end{Highlighting}
\end{Shaded}

Many teams use tools like \textbf{Renovate} or \textbf{Dependabot} to
automatically open PRs when a newer base image is available.

\subsection{Putting It All Together}\label{putting-it-all-together}

Here is a Dockerfile for our Nginx image that applies all the practices
above:

\begin{Shaded}
\begin{Highlighting}[]
\KeywordTok{FROM}\NormalTok{ nginx:alpine}

\CommentTok{\# Copy content}
\KeywordTok{COPY} \OperatorTok{{-}{-}chown=nginx:nginx}\NormalTok{ html /usr/share/nginx/html}

\CommentTok{\# Drop to non{-}root}
\KeywordTok{USER}\NormalTok{ nginx}

\KeywordTok{EXPOSE}\NormalTok{ 80}

\KeywordTok{CMD}\NormalTok{ [}\StringTok{"nginx"}\NormalTok{, }\StringTok{"{-}g"}\NormalTok{, }\StringTok{"daemon off;"}\NormalTok{]}
\end{Highlighting}
\end{Shaded}

And the corresponding \texttt{docker-compose.yml}:

\begin{Shaded}
\begin{Highlighting}[]
\FunctionTok{services}\KeywordTok{:}
\AttributeTok{  }\FunctionTok{frontend}\KeywordTok{:}
\AttributeTok{    }\FunctionTok{build}\KeywordTok{:}\AttributeTok{ .}
\AttributeTok{    }\FunctionTok{image}\KeywordTok{:}\AttributeTok{ my{-}nginx:secure}
\AttributeTok{    }\FunctionTok{container\_name}\KeywordTok{:}\AttributeTok{ frontend}
\AttributeTok{    }\FunctionTok{ports}\KeywordTok{:}
\AttributeTok{      }\KeywordTok{{-}}\AttributeTok{ }\StringTok{"8080:80"}
\AttributeTok{    }\FunctionTok{read\_only}\KeywordTok{:}\AttributeTok{ }\CharTok{true}
\AttributeTok{    }\FunctionTok{tmpfs}\KeywordTok{:}
\AttributeTok{      }\KeywordTok{{-}}\AttributeTok{ /tmp}
\AttributeTok{      }\KeywordTok{{-}}\AttributeTok{ /var/cache/nginx}
\AttributeTok{      }\KeywordTok{{-}}\AttributeTok{ /var/run}
\AttributeTok{    }\FunctionTok{cap\_drop}\KeywordTok{:}
\AttributeTok{      }\KeywordTok{{-}}\AttributeTok{ ALL}
\AttributeTok{    }\FunctionTok{cap\_add}\KeywordTok{:}
\AttributeTok{      }\KeywordTok{{-}}\AttributeTok{ NET\_BIND\_SERVICE}
\AttributeTok{    }\FunctionTok{networks}\KeywordTok{:}
\AttributeTok{      }\KeywordTok{{-}}\AttributeTok{ app{-}network}

\FunctionTok{networks}\KeywordTok{:}
\AttributeTok{  }\FunctionTok{app{-}network}\KeywordTok{:}
\AttributeTok{    }\FunctionTok{driver}\KeywordTok{:}\AttributeTok{ bridge}
\end{Highlighting}
\end{Shaded}

\subsection{Conclusion}\label{conclusion-6}

In this article we covered:

\begin{itemize}
\tightlist
\item
  Running containers as a \textbf{non-root user} with \texttt{USER} in
  the Dockerfile
\item
  Using a \textbf{read-only filesystem} with \texttt{-\/-read-only} and
  \texttt{tmpfs} for writable paths
\item
  \textbf{Dropping Linux capabilities} with \texttt{-\/-cap-drop=ALL}
  and adding back only what is needed
\item
  Keeping \textbf{secrets out of images and environment variables}
\item
  Choosing \textbf{minimal base images} to reduce attack surface
\item
  \textbf{Scanning images} for vulnerabilities with Docker Scout or
  Trivy
\item
  \textbf{Keeping images up to date} with pinned digests and automated
  updates
\end{itemize}

This concludes the \textbf{Getting Started with Docker} series. You now
have a solid foundation to build, network, persist, orchestrate, and
secure containerized applications. The natural next step is
\textbf{Docker Swarm} for multi-host deployments, or \textbf{Kubernetes}
for large-scale container orchestration.

\begin{center}\rule{0.5\linewidth}{0.5pt}\end{center}

If you enjoyed this article, don't forget to \textbf{give it a clap 👏},
\textbf{share it with your friends 🔗}, and \textbf{follow me for more
tips and tutorials on software development 📘}. Your support helps me
create more content like this --- thank you! 🙌

\end{document}
